% ABNT - PACOTE DE ACENTUAÇÃO
\usepackage[utf8]{inputenc} 

% MARGENS LEFT RIGHT TOP E BOTTOM
\usepackage[lmargin=2cm,
            tmargin=0.6cm,
            rmargin=2cm,
            bmargin=3cm,
            headsep=0.0cm,
            includeheadfoot]{geometry}

\setlength\headheight{2.8cm}

% LEVA O ESTILO DO TEXTO PARA OUTROS PROGRAMAS TIPO WORD
\usepackage[T1]{fontenc}

% PÕE O PRIMEIRO PARÁGRAFO AUTOMATICAMENTE
\usepackage{indentfirst}

% HABILITA A UTILIZAÇÃO DE FONTES
\renewcommand{\familydefault}{\sfdefault}

\usepackage[onehalfspacing]{setspace}  % ESPAÇAMENTO ENTRE AS LINHAS
\usepackage[brazil]{babel}             % TRADUZ PARA PORTUGUÊS
\usepackage{comment}                   % COLOCAR COMENTÁRIOS
\usepackage{enumerate}                 % PACOTE PARA ENUMERAÇÃO
\usepackage{multirow, multicol}        % FORMATAR TABELAS
\usepackage{colortbl}                  % TABULAR COM COR 
\usepackage{lipsum}                    % TEXTO ALEATORIO
\usepackage{dashrule}                  % USAR LINHAS PONTILHADAS
\usepackage{paracol}                   % TEXTOS EM COLUNAS
\usepackage{hyperref}                  % HIPERLINKS NO SUMÁRIO

\pagestyle{empty}       % PREENCHIMENTO DA PÁGINA (FANCY HABILITA)
\usepackage{fancyhdr}   % CABEÇALHO E RODAPÉ
\fancyhf{}              % SEM NUMERAÇÃO NO MEIO PÁGINA

\usepackage{graphicx}   % PACOTE PARA IMAGENS
\usepackage{efbox}      % PACOTE PARA BORDAS EM IMAGENS
\usepackage{subfig}     % IMAGENS LADO A LADO
\usepackage{float}      % POSICIONAR IMAGEM NO COMANDO FIGURE
\usepackage{tabularx}   % PARA TABULAR MAIS BONITO

% HABILITAR CORES PARA TEXTOS E CAIXAS DE TEXTO COLORIDAS
\usepackage[x11names]{xcolor}

% COR TOPICOS \SECTION{}
\usepackage{sectsty}

% CAIXA DE TEXTO COLORIDA COM TÍTULO
\usepackage{wrapfig}
\usepackage[most]{tcolorbox}

% MATEMÁTICA
\usepackage{mathrsfs}
\usepackage{amsmath}
\usepackage{amsthm}
\usepackage{amsfonts}
\usepackage{amssymb}
\usepackage{dsfont}
\usepackage{mathtools}
\usepackage{tikz}                        
\usetikzlibrary{babel,quotes, angles}   % ÂNGULOS
\usepackage{esvect}                     % VETORES \vv{}
\everymath{\displaystyle}               % FRAÇÕES MAIS BONITAS

% \cancel{x}, \bcancel{x} (invertido), \cancelto{x}{para y}   
\usepackage{cancel}

% OUTROS
%\usepackage{parskip}        % ALINHAMENTO TEXTO (ELIMINA PARÁGRAFOS)
\usepackage{blindtext}      % IMPRIME TEXTO ALEATÓRIO \blindtext
\usepackage{fancybox}       % MARGENS CUSTOMIZÁVEIS

% FONTE MATEMÁTICA
\DeclareMathVersion{testemath}
\SetSymbolFont{operators}{testemath}{OT1}{bch}{n}{n}
\SetSymbolFont{letters}{testemath}{OML}{bch}{n}{n}
\mathversion{testemath}